% !TEX TS-program = pdflatex
% !TeX program = pdflatex
% !TEX encoding = UTF-8
% !TEX spellcheck = en_US

\documentclass{beamer}


\def\accyear{2024-2025}

\institute{
	Laboratoire de la Communication dans les Systèmes Informatiques (LCSI)\\
	École  nationale Supérieure d'Informatique (ESI, ex. INI), Algiers, Algeria
}

\author[Abdelkrime Aries (\accyear)] %
{Abdelkrime Aries}

\titlegraphic{%
	\includegraphics[height=1cm]{../extra/logo/esi-logo.png}%
	\hspace*{2cm}%
	\includegraphics[height=1cm]{../extra/logo/lcsi-logo.png}%
	\hspace*{2cm}%
	\includegraphics[height=1cm]{../extra/logo/esi.nlp.pdf}%
}%\hspace*{4.75cm}~

\date{Academic year: \accyear} %\today

\title[NLP]{Natural Language Processing} 

\usefolder{../extra/beamer}

\usetheme{Karimnlp} % Antibes Boadilla Warsaw

\beamertemplatenavigationsymbolsempty



%\usepackage{stix}

\title[ESI - NLP\_master: Presentation] %
{Natural language processing (Master)\\Course (also, curse) presentation} 

\setbeamertemplate{headline}{}

\begin{document}
 
 
 \begin{frame}
 	\frametitle{Course resources}
 	
 	\begin{itemize}
 		\item Course project
 		\begin{itemize}
 			\item English: \url{https://github.com/projeduc/ESI_NLP}
 			\item French: \url{https://github.com/projeduc/ESI_TALN}
 		\end{itemize}
 		\item Project content
 		\begin{itemize}
 			\item demos: demonstrations of some well known NLP tools 
 			\item labs: previous labs (graded practicals)
 			\item workshops: previous workshops (ungraded practicals)
 			\item tutos: tutorials (exercises) 
 			\item pres: presentations
 		\end{itemize}
 	\end{itemize}
 	
 \end{frame}

\begin{frame}
\frametitle{Prerequisites}

\begin{itemize}
	\item Basic knowledge of languages (vocabulary, grammar, etc.);
	\item Data structures (1CP-ALSDD);
	\item Probability and statistics (2CP-PRST1, 2CP-PRST2);
	\item Algebra and calculus (2CP-ANAL3, 2CP-ALG3, 2CP-ANAL4);
	\item Language theory (1CS-THP);
	\item Operations research (1CS-RO);
	\item Machine learning (2CS-ML or 2CS-BDM);
	\item Programming in Python.
\end{itemize}

\end{frame}


\begin{frame}
\frametitle{Objectives}

\begin{itemize}
	\item Learn some linguistic concepts;
	\item Apply concepts covered in the prerequisite modules to solve NLP tasks;
	\item Discover some NLP tools and resources;
	\item Understand the research directions in this field;
	\item Acquire the ability to articulate a novel NLP methodology and assess its viability within the framework of a scholarly publication.
\end{itemize}

\end{frame}


\begin{frame}
\frametitle{Course components \& Assessment}

\optword{Course components:}
\begin{itemize}
	\item \textbf{Lectures:} 15 hours.
	\item \textbf{Quizzes:} online, the end of the day before midnight (graded).
	\item \textbf{Labs:} 9 hours (graded).
	\item \textbf{Online Demos:} {\scriptsize\url{https://github.com/projeduc/ESI_NLP/tree/main/demos}}.
\end{itemize}

\vfill

\optword{Assessment:}
\begin{itemize}
	\item \textbf{Quizzes:} 20\%.
	\item \textbf{Labs:} 30\%.
	\item \textbf{Written exam:} 50\%.
\end{itemize}

\end{frame}

\begin{frame}
\frametitle{Lectures}

\begin{enumerate}
	\item Introduction to NLP;
	\item ML for NLP;
	\item Basic text processing;
	\item Language models;
	\item Part of Speech (PoS) tagging;
	\item Word semantics;
	\item Some applications.
\end{enumerate}

\end{frame}

\begin{frame}
	\frametitle{Quizzes}
	
	\begin{itemize}
		\item Answering some questions about the lecture.
		\item \optword{GOAL:} Encouraging course review.
		\item \optword{Method:} Take a Quiz (offline).
		\begin{itemize}
			\item \textbf{When?} Each end of day.
			\item \textbf{What?} Google Form.
			\item \textbf{How much?} 3 Quizzes.
			\item \textbf{Duration?} 20 minutes/quiz max.
			\item \textbf{Deadline?} midnight (the same day).
		\end{itemize}
		\item \optword{Assessment:}
		\begin{itemize}
			\item 20 points (the mark is an integer).
			\item In a multiple-choice question, a wrong choice cancels a right one.
			\item For each question, the min mark is 0 (no less).
		\end{itemize}
	\end{itemize}
	
\end{frame}

\begin{frame}
	\frametitle{Labs}
	
	\begin{itemize}
		\item \textbf{Lab01}: Discovery of some NLP tools.
		\begin{itemize}
			\item \optword{GOAL:} Learn how to use some NLP tools.
			\item \optword{Method:} Follow instructions to complete an application.
		\end{itemize}
		\item \textbf{Lab02}: Developing an NLP application.
		\begin{itemize}
			\item \optword{GOAL:} Learn how to enhance an existing application.
			\item \optword{Method:} Modify an existing application to improve it based on some metrics.
		\end{itemize}
		\item \textbf{Lab03}: Writing an article on this application.
		\begin{itemize}
			\item \optword{GOAL:} Learn \LaTeX\ and NLP-related critical writing.
			\item \optword{Method:} Write a paper about the new method (Lab02).
		\end{itemize}
	\end{itemize}
	
\end{frame}

\begin{frame}
	\frametitle{Color codes}
	
	
	\begin{itemize}
		\item Slides' texts and blocks
		\begin{itemize}
			\item {\Huge${\color{red}\blacksquare}$} keywords
			\item {\Huge${\color{violet}\blacksquare}$} options (choices)
			\item {\Huge${\color{olivegreen}\blacksquare}$} examples
		\end{itemize}
		\item Figures (mostly)
		\begin{itemize}
			\item {\Huge${\color{red}\blacksquare}$} process
			\item {\Huge${\color{blue}\blacksquare}$} inputs
			\item {\Huge${\color{green}\blacksquare}$} intermediate states
		\end{itemize}
	\end{itemize}
	
\end{frame}

\begin{frame}
	\frametitle{Similar courses}
	
	\begin{itemize}
			
			\item Natural Language Processing S21 (updated).
			Carnegie Mellon University. 
			\kurl{http://demo.clab.cs.cmu.edu/NLP/}
			
			\item CS388: Natural Language Processing (2018). 
			University of Texas. 
			\kurl{https://www.cs.utexas.edu/~mooney/cs388/}
			
			\item CSE 517: Natural Language Processing ().
			University of Washington.
			\kurl{https://courses.cs.washington.edu/courses/cse517/}
			
			\item CS 294-5: Statistical Natural Language Processing (2005). 
			University of Berkely. 
			\kurl{https://people.eecs.berkeley.edu/~klein/cs294-5/index.html}
			
			\item CS224n: Natural Language Processing with Deep Learning (updated).
			Stanford university.
			\kurl{https://web.stanford.edu/class/cs224n/}
			
			\item CS 5340/6340: Natural Language Processing (updated). 
			University of Utah.
			\kurl{https://my.eng.utah.edu/~cs5340/schedule.html}
		
	\end{itemize}
	
\end{frame}

\begin{frame}
	\frametitle{Reading resources}
	
	\begin{itemize}
		
		\item \textbf{Traitement automatique du langage naturel - course textbook [FR]}
		
		by Abdelkrime Aries
		
		\kurl{https://github.com/projeduc/ESI_2CS_TALN/releases/download/4.0/taln_v2.pdf}
		
		
		\item \textbf{Speech and Language Processing} 
		
		by Dan Jurafsky and James H. Martin
		
		\kurl{https://web.stanford.edu/~jurafsky/slp3/}
		
		\item \textbf{Natural language processing}
		
		by Jacob Esenstein 
		
		\kurl{https://cseweb.ucsd.edu/~nnakashole/teaching/eisenstein-nov18.pdf}
		
		\item \textbf{Modern Approaches in Natural Language Processing - course textbook (LMU Munich)}
		
		by Carolin Becker, et al.
		
		\kurl{https://slds-lmu.github.io/seminar_nlp_ss20/index.html}
		
		
	\end{itemize}
	
\end{frame}

\begin{frame}
	\frametitle{Disclaimer}
	
	\begin{tcolorbox}[colback=red,
		colframe=red,  
		arc=0pt,outer arc=0pt,
		valign=top, 
		halign=center,
		width=\textwidth]
		
		\color{white}
		{\bfseries Disclaimer}
		
	\end{tcolorbox}\vspace{-.5cm}
	\begin{tcolorbox}[colback=my-grey,
		colframe=my-grey,  
		center, arc=0pt,outer arc=0pt,
		valign=top, 
		halign=left,
		width=\textwidth]
		
		\begin{itemize}
			\item An engineer must be a force for good; we must not use our knowledge to harm others.
			\item This course does not encourage malicious use of NLP techniques such as espionage, illegal information gathering, misinformation via chatbots, etc.
			\item Any example that may contain one of these applications is purely educational.	
		\end{itemize}
		
	\end{tcolorbox}
	
\end{frame}


\end{document}

